\documentclass[11pt,a4paper]{article}

\usepackage[T1]{fontenc}
\usepackage[utf8]{inputenc}
\usepackage{lmodern}
\usepackage{microtype}
\usepackage{amsmath}
\usepackage{amssymb}
\usepackage[margin=2.6cm]{geometry}
\usepackage{booktabs}
\usepackage{tabularx}
\usepackage{graphicx}
\usepackage{caption}
\usepackage{float}
\usepackage{xcolor}
\usepackage{fancyhdr}
\usepackage[hidelinks]{hyperref}

\definecolor{accent}{HTML}{275CB2}
\definecolor{notebg}{HTML}{F1F5FB}
\hypersetup{colorlinks=true, linkcolor=accent, urlcolor=accent}
\captionsetup{font=small, labelfont={bf,color=accent}, width=0.92\textwidth}

\pagestyle{fancy}
\fancyhf{}
\renewcommand{\headrulewidth}{0.4pt}
\fancyhead[L]{\footnotesize\sffamily Project 4 --- Preference Elicitation}
\fancyhead[R]{\footnotesize\sffamily Kartik Anand, 676049}
\fancyfoot[C]{\footnotesize\thepage}

\setlength{\parskip}{0.55em}
\setlength{\parindent}{0pt}
\renewcommand{\arraystretch}{1.15}
\setlength{\emergencystretch}{2.5em}

\makeatletter
\renewcommand\section{\@startsection{section}{1}{\z@}%
  {-3.5ex \@plus -1ex \@minus -.2ex}{2.2ex \@plus.2ex}%
  {\normalfont\Large\sffamily\bfseries\color{accent}}}
\renewcommand\subsection{\@startsection{subsection}{2}{\z@}%
  {-3.0ex \@plus -1ex \@minus -.2ex}{1.4ex \@plus .2ex}%
  {\normalfont\large\sffamily\bfseries\color{accent}}}
\renewcommand\paragraph{\@startsection{paragraph}{4}{\z@}%
  {2.4ex \@plus1ex \@minus.2ex}{-0.8em}%
  {\normalfont\normalsize\sffamily\bfseries\color{black}}}
\makeatother

\newsavebox{\calloutbuf}
\newenvironment{callout}
  {\par\medskip\begin{lrbox}{\calloutbuf}%
   \begin{minipage}{\dimexpr\linewidth-24pt\relax}\setlength{\parskip}{0.4em}}
  {\end{minipage}\end{lrbox}%
   \noindent\colorbox{notebg}{\hspace{4pt}\usebox{\calloutbuf}\hspace{4pt}}\par\medskip}

\newcommand{\file}[1]{\texttt{\small #1}}

% enumitem is absent from this TeX installation, so the two list shapes it was
% supplying are built from the base classes: a tightened enumerate, and a
% description whose label sits on its own line.
\newenvironment{steps}
  {\begin{enumerate}\setlength{\itemsep}{0.15em}\setlength{\parskip}{0.15em}}
  {\end{enumerate}}
\newcommand{\term}[1]{\par\medskip\noindent\textbf{#1}\par\nopagebreak\vspace{0.15em}}
\newenvironment{terms}{\par\begingroup\leftskip=1.2em}{\par\endgroup\medskip}

\title{\vspace{-1.4cm}\sffamily\color{accent}\bfseries
       Preference Elicitation for a Cold-Start Movie Recommender\\[0.25em]
       \Large\mdseries Project 4 --- Human-Centric Artificial Intelligence}
\author{Kartik Anand \quad---\quad Matriculation 676049\\[0.3em]
        \small Hamburg University of Technology \quad SoSe 2026}
\date{\small \today}

\begin{document}
\maketitle
\thispagestyle{fancy}

\begin{callout}
A new user has rated nothing, so a recommender must ask. This report chooses a
20-dimensional movie representation and argues that the choice is governed by the
elicitation budget rather than by what describes a film best; extends
Bradley--Terry to full rankings via Plackett--Luce and shows why that particular
extension is the one that makes the study's comparison meaningful; and designs a
150-participant, time-matched, pre-registered study comparing pairwise choice
against ranking ten. A power simulation overturns the design's own starting
intuition: one ranking of ten carries about sixty times the Fisher information of
one pairwise comparison, and yet at an equal time budget the pairwise interface
recovers $w$ at least as well, because a ranking's information is concentrated on
the ten films in front of the participant rather than spread across the feature
space. The study is designed, not run. The interface that would run it is
implemented and reachable from the project landing page.
\end{callout}

\tableofcontents
\vspace{1em}

% ============================================================================
\section{Setting}

The IMDB 5000 dataset describes roughly 5{,}000 films and contains no user
ratings at all. That absence is not a defect to work around; it is the situation
of every genuinely new user of every recommender, and the only way out of it is
to ask.

The preference model is fixed by the brief. A film is a feature vector
$x \in \mathbb{R}^d$, a user is a latent weight vector $w \in \mathbb{R}^d$, and
utility is linear:
\[
  U(x) \;=\; w^{\!\top} x .
\]
Elicitation is the business of recovering $w$ from a small number of
interactions. Two interfaces are on the table --- repeated choice between two
films, and ranking sets of ten --- and the project is to work out which of them
buys more information per unit of the participant's patience.

Everything is numpy and Django. The likelihood, its gradient, the optimiser and
the Fisher information are written out in \file{project4/preference.py} rather
than imported, because the derivation is the deliverable.

% ============================================================================
\section{Task 1 --- The feature representation}

\subsection{The constraint that actually decides it}

The instinct is to ask which columns best describe a film. That is the wrong
question here, and answering it produces a representation that cannot be
estimated. Two constraints decide the feature set, and neither is about
descriptive power.

\paragraph{The elicitation budget.} $w$ will be fitted from something like
twenty-five responses. Consider adding one weight per director. The dataset has
2{,}398 distinct directors; two films drawn uniformly at random share one with
probability well under $10^{-3}$. Essentially every response a participant will
ever give multiplies every one of those 2{,}398 weights by zero. They are not
weak features --- they are \emph{unidentifiable} ones at this sample size, and
including them does not add information, it adds room for the prior to invent
preferences that were never expressed. The same argument disposes of the three
actor-name columns.

\paragraph{Visibility to the participant.} A feature the participant cannot see
on the card cannot have caused their choice. \file{budget}, \file{gross}, the
four Facebook-likes columns, \file{aspect\_ratio} and
\file{facenumber\_in\_poster} are not consulted by anyone picking a film to watch
tonight. If such a feature sits in $x$ but not on the card, the fit will
nonetheless attribute choices to it whenever it happens to correlate with
something that did matter. The rule adopted is symmetric and easy to check:
everything in $x$ appears on the card, and everything on the card is in $x$.

\subsection{Which genres earn a dimension}

Genres are pipe-separated in the source and 26 distinct tags appear. Rather than
picking a round number, the cut follows from the budget argument. Draw two films
uniformly at random. For a genre present in a fraction $p$ of the pool, they
differ on it with probability
\[
  \Pr[\text{informative on genre } g] \;=\; 2p(1-p),
\]
and a pair that agrees on a genre carries no information about its weight at all.
Requiring at least one informative pair in ten gives $2p(1-p) \geq 0.10$, whose
lower root is $p \geq 0.0528$. Fourteen genres clear it.

\begin{table}[H]
\centering\small
\begin{tabularx}{\textwidth}{lX}
\toprule
\textbf{Kept} & Drama, Comedy, Thriller, Action, Adventure, Romance, Crime,
                Sci-Fi, Fantasy, Mystery, Horror, Family, Biography, Animation \\
\textbf{Dropped} & War, History, Sport, Music, Musical, Western, Documentary,
                Film-Noir, News \\
\bottomrule
\end{tabularx}
\caption{The genre cut at $p \geq 0.0528$. Western sits at $p = 0.019$, so about
one pair in twenty-six says anything about it --- one observation across a whole
session, which is not enough to estimate a weight and is enough for the prior to
fill one in.}
\end{table}

\subsection{The remaining six, and their scaling}

Four continuous axes and two indicators complete the vector, giving $d = 20$:
how recent the film is (\file{title\_year}), how long it runs
(\file{duration}), how well reviewed it is (\file{imdb\_score}), how widely seen
(\,$\log_{10}$ of \file{num\_voted\_users}\,), whether it is certified for adults,
and whether it is black and white.

Acclaim and reach are deliberately separate dimensions rather than one
``popularity'' axis. ``Well reviewed but obscure'' and ``widely seen but
mediocre'' are both common and describe genuinely different tastes; collapsing
them would make those two people indistinguishable.

The four continuous features are standardised, so $w_j$ reads uniformly as
utility per standard deviation and a single isotropic prior is a coherent thing
to put on them. The genre indicators are deliberately \emph{not} standardised:
scaling a genre with $p = 0.06$ to unit variance would let the same weight
produce a four-times-larger utility swing than for a genre with $p = 0.5$, so an
isotropic prior would permit rarer genres --- the ones with less evidence behind
them --- larger effects. Leaving them as raw $0/1$ makes $w_g$ the utility of the
tag itself and shrinks every genre equally.

\subsection{The pool}

Films with fewer than 25{,}000 IMDB votes are removed and duplicate titles
collapsed, leaving \textbf{2{,}734} films spanning 1927--2016. This is a validity
decision, not data cleaning: a participant comparing two films they have never
heard of is not expressing a preference, they are reading metadata off the card
and guessing. That guessing enters the model as response noise attributed to
$w$, and no amount of modelling recovers from it.

Extraction is \file{project4/data.py}, function \file{pool()}. The threshold is
computed from \file{PAIR\_INFORMATIVENESS} rather than hard-coded, so the
argument above and the code cannot drift apart.

% ============================================================================
\section{Task 2 --- From Bradley--Terry to rankings}

\subsection{The extension}

Bradley--Terry over linear utilities is
\[
  \Pr(i \succ j) \;=\; \frac{e^{U_i}}{e^{U_i} + e^{U_j}}
                 \;=\; \sigma\!\left(w^{\!\top}(x_i - x_j)\right).
\]

Read a ranking as a sequence of choices rather than as an object in its own
right. The participant names a favourite out of ten; then a favourite of the nine
that remain; and so on. Each of those is Bradley--Terry widened from two
alternatives to however many are left, which is a softmax over the surviving
utilities. Multiplying the stages gives, for the observed order
$i_1 \succ i_2 \succ \cdots \succ i_n$,
\begin{equation}
  \Pr(i_1 \succ \cdots \succ i_n)
  \;=\; \prod_{k=1}^{n-1} \frac{e^{U_{i_k}}}{\sum_{l \geq k} e^{U_{i_l}}} ,
  \label{eq:pl}
\end{equation}
the Plackett--Luce model. The final factor is $1$ and is dropped.

\subsection{Why this extension and not another}

\paragraph{It is not a replacement.} At $n = 2$ the product in \eqref{eq:pl} has
a single factor and reduces to Bradley--Terry exactly. The two interfaces in this
study are therefore not fitting two different models, and \file{preference.py}
never branches on which arm produced a response --- a pairwise choice is stored
as a ranking of length two and the same likelihood consumes both.

\paragraph{Its pairwise marginals are Bradley--Terry.} This is the property the
entire study rests on. Under \eqref{eq:pl}, the probability that $i$ appears
anywhere before $j$ in the sampled ranking is
\[
  \Pr(i \text{ before } j) \;=\; \frac{e^{U_i}}{e^{U_i}+e^{U_j}},
\]
independently of which other films share the set. So an answer collected by
ranking ten films and an answer collected by comparing two carry information
about the \emph{same} parameter. Without this, scoring both arms against a common
set of held-out pairwise comparisons (Section~\ref{sec:outcome}) would not be a
fair test --- it would be a category error, measuring one arm on its own terms
and the other on somebody else's. The claim is verified numerically rather than
asserted: \file{check\_marginals} draws $6\times10^4$ Plackett--Luce rankings by
Gumbel-max sampling and compares the empirical order frequencies against the
Bradley--Terry probabilities, agreeing to a maximum absolute error of
$3.3\times10^{-3}$, which is Monte-Carlo noise at that sample size.

\paragraph{The tempting alternative is wrong.} A ranking of ten implies
$\binom{10}{2} = 45$ pairwise preferences, and it is natural to treat them as 45
Bradley--Terry observations and multiply. This fails twice. The 45 comparisons
are not independent: $i_1 \succ i_3$ is partly entailed by $i_1 \succ i_2$ and
$i_2 \succ i_3$, so the product counts the same evidence repeatedly and reports
far more confidence than the data support --- exactly the failure that would make
Design 2 look artificially good. And the product is not a probability
distribution over the $n!$ orderings; it does not sum to one, so it is not a
likelihood and the estimate that maximises it is not a maximum-likelihood
estimate of anything. Equation~\eqref{eq:pl} is a product of properly normalised
choices and sums to one by construction.

\paragraph{What it inherits, including the awkward part.} Plackett--Luce carries
over Luce's choice axiom: the odds between two films do not depend on what else
is present. People violate this. Put two near-identical sequels into a set of ten
and they tend to split the appeal and both sink. Design 2 exposes ten films at
once and is therefore more exposed to the violation than Design 1, which never
shows more than two. This is recorded as a limitation and as a threat to validity
in Section~\ref{sec:threats}, not quietly assumed away.

\subsection{Estimating $w$}

The gradient of \eqref{eq:pl} is clean. At stage $k$ it is the chosen film's
features minus the softmax-weighted mean over the films still available:
\[
  \nabla_w \log \Pr \;=\; \sum_{k=1}^{n-1}
    \Bigl( x_{i_k} - \sum_{l \geq k} p^{(k)}_l\, x_{i_l} \Bigr),
  \qquad
  p^{(k)}_l = \frac{e^{U_{i_l}}}{\sum_{m \geq k} e^{U_{i_m}}} .
\]

The estimate is MAP rather than ML, with an isotropic Gaussian prior
$w \sim \mathcal{N}(0, \sigma^2 I)$ at $\sigma = 1$. The prior is not a
refinement; without it the estimate frequently does not exist. With $d = 20$ and
roughly 25 responses there are usually directions in $w$ that no response
constrains, and the likelihood increases without bound along them. $\sigma = 1$
states that a one-standard-deviation move in a feature is worth about one unit of
utility --- roughly a $73/27$ preference --- which is a weak claim and the honest
one: absent evidence, assume no strong opinion.

Optimisation is gradient ascent with backtracking line search. The objective is a
concave log-likelihood plus a strictly concave penalty, hence strictly concave, so
a vanishing gradient certifies the global optimum; \file{fit} reports the final
gradient norm and the interface displays it. Typical runs terminate around
$4\times10^{-7}$.

% ============================================================================
\section{Task 3 --- The user study}

\subsection{What is being compared, and the arithmetic that motivates it}

\paragraph{Information per task.} For a softmax choice over a set $S$ with
probabilities $p_s$, the Fisher information about $w$ is
$\sum_{s \in S} p_s x_s x_s^{\!\top} - \bar{x}\bar{x}^{\!\top}$ with
$\bar{x} = \sum_s p_s x_s$, and a Plackett--Luce ranking contributes one such
term per stage. Evaluated on this feature set at a representative $w$, the trace
is about $0.57$ for one pairwise comparison and about $34.1$ for one ranking of
ten --- a factor of roughly \textbf{60}. Taken alone this looks decisive for
Design 2.

\paragraph{The confound that would ruin the comparison.} Comparing ``forty pairs
against eight rankings'' compares two arbitrary numbers. What a recommender
actually spends is the user's patience, and eight minutes is eight minutes
however it is divided. \textbf{Both arms therefore receive the same wall-clock
elicitation budget of eight minutes} and complete as many tasks as fit. This is
the single most important control in the design, and applying it immediately
deflates the factor of 60: at the pilot's assumed 6\,s per comparison and 90\,s
per ranking, the budget buys 80 comparisons or 5 rankings, so the accumulated
traces come out at 118 and 129 --- a ratio of \textbf{1.09}, not 60.

\paragraph{And the trace is the wrong summary anyway.}
\label{sec:conditioning}
Trace measures total information, not whether it is pointed anywhere useful.
Accumulating the Fisher matrix over a whole session and looking at its spectrum
tells a different story.

\begin{table}[H]
\centering\small
\begin{tabular}{lcccc}
\toprule
& \textbf{Trace} & \textbf{Log-determinant} & \textbf{Smallest eigenvalue} &
\textbf{Distinct films} \\
\midrule
Design 1 --- 80 pairs    & 118.3 & $+17.47$ & 0.224 & 160 \\
Design 2 --- 5 rank-tens & 129.2 & $-0.09$  & 0.035 & \phantom{0}50 \\
\bottomrule
\end{tabular}
\caption{Fisher information accumulated over one eight-minute session, averaged
over 40 draws of $w$. The traces nearly tie; the determinants do not.}
\end{table}

The log-determinant --- the D-optimality criterion, and the quantity that
actually governs the volume of the confidence region for $\hat{w}$ --- is
\textbf{17.6 nats lower} for Design 2, and its smallest eigenvalue is
\textbf{6.4 times smaller}. The reason is structural rather than statistical.
Eighty pairwise comparisons touch 160 films and contribute 160 independent random
directions in feature space. Five rankings touch 50 films, and the 45 stage-wise
comparisons inside one ranking all lie in the span of that set's ten feature
vectors. A ranking piles up a great deal of information about a few directions
and leaves others essentially unconstrained, where the prior then supplies the
answer.

This is the substantive contribution of the analysis, and it is easy to miss: the
per-task factor of 60 is real, the time correction cancels almost all of it, and
what remains points the other way.

\paragraph{One more reason to doubt Design 2.} The arithmetic above assumes the
participant answers all nine stages of a ranking exactly as the model says.
Ordering the middle of a list of ten is genuinely hard, and the natural failure
mode is position-dependent --- careful at the top, careless by position eight. If
the later stages are noise, even Design 2's remaining advantage evaporates.

\subsection{Hypotheses}

Stated before data collection, with one confirmatory primary and three
subordinate secondaries.

\begin{terms}
\term{H1 (primary, confirmatory).} Under an equal eight-minute budget, $w$
  estimated from Design 1 (pairwise) predicts held-out pairwise choices at least
  as well as $w$ estimated from Design 2 (rank ten), measured by mean held-out
  log loss. Directional, one-sided at $\alpha = .05$, in favour of Design 1.

  \emph{The direction is worth pausing on.} It is the opposite of what the
  per-task information figure suggests, and it was set by the analysis in
  Section~\ref{sec:conditioning} and the simulation in Section~\ref{sec:sim} ---
  both run before this hypothesis was written down, and both reported here
  whichever way they came out. Pre-registering the direction that the naive
  arithmetic favours, and then finding the opposite, is precisely the situation
  pre-registration exists to make visible.
\term{H2 (secondary).} Response quality in Design 2 degrades with position: the
  implied choice at stage $k$ is noisier as $k$ grows.
\term{H3 (secondary).} Design 2 has a higher abandonment rate and lower
  self-reported ease.
\term{H4 (secondary).} Design 1's advantage narrows as the budget grows. Once
  enough rankings have accumulated to cover the feature space, the conditioning
  problem in Section~\ref{sec:conditioning} eases, and Design 2's larger per-task
  information should begin to tell.
\end{terms}

\subsection{The primary outcome, and why it cannot be ``accuracy of $w$''}
\label{sec:outcome}

The obvious outcome measure is how close the estimated $\hat{w}$ is to the
participant's true $w$. It is unavailable: $w$ is latent and never observed for a
real person, so no distance to it can be computed. This is not a technicality to
be finessed --- a study whose primary outcome cannot be computed is not a study.

What is measurable is prediction. Every participant, in both arms, finishes with
an identical block of \textbf{ten held-out pairwise comparisons}. $\hat{w}$ is
fitted on the elicitation responses only, and scored on the held-out block by
mean log loss
\[
  \mathcal{L} \;=\; -\frac{1}{10}\sum_{m=1}^{10} \log
    \sigma\!\left(w^{\!\top}(x_{c_m} - x_{r_m})\right),
\]
where $c_m$ was chosen and $r_m$ rejected. An uninformed model answers $0.5$
every time and scores $\log 2 = 0.693$; anything below has learned something.

Log loss rather than accuracy, for a concrete reason: on ten held-out items,
accuracy takes eleven possible values and discards all confidence information, so
it needs far more participants to separate the arms. The held-out block is
identical in task, count and wording across arms --- if the test itself differed,
any difference in the result could be the test.

\subsection{Design and assignment}

Between-subjects, with a within-subject held-out block. Each participant does one
elicitation design, assigned at random.

Within-subjects would be more sensitive, and is rejected for a reason that no
counterbalancing removes: having ranked ten films, a participant understands
their own taste better and would approach the second design differently for that
reason alone. Order effects can be balanced; genuine learning about oneself
cannot be undone.

Randomisation is blocked on the pre-task question ``how often do you watch a
film'' (four levels), in permuted blocks of four within each stratum. Frequent
viewers have firmer preferences and lower response noise, and leaving that to
simple randomisation would let it land unevenly across arms.

\begin{table}[H]
\centering\small
\begin{tabularx}{\textwidth}{lXX}
\toprule
& \textbf{Design 1} & \textbf{Design 2} \\
\midrule
Elicitation task & choose one of two & rank ten, most to least preferred \\
Budget & 8 minutes & 8 minutes \\
Expected tasks completed & $\approx 80$ at 6\,s each & $\approx 5$ at 90\,s each \\
Model & Bradley--Terry & Plackett--Luce, Eq.~\eqref{eq:pl} \\
Held-out block & 10 pairs, identical & 10 pairs, identical \\
\bottomrule
\end{tabularx}
\caption{The two arms. Per-task durations are the pilot's job to confirm; the
budget, not the task count, is what is held fixed.}
\end{table}

Films are drawn uniformly at random from the pool, as the brief permits. Adaptive
selection would raise both arms' efficiency but is deliberately excluded: it
would interact with the design being tested, since an adaptive scheme has more
room to exploit a ranking of ten than a single pair, and the comparison would no
longer isolate the interface.

\subsection{Procedure}

\begin{steps}
\item Consent form: what is recorded, what is not, right to withdraw.
\item Viewing-frequency question; blocked random assignment.
\item Instructions with one worked practice task, not counted.
\item Eight minutes of elicitation in the assigned design. A visible timer. The
  participant may stop early; responses so far are retained.
\item Ten held-out pairwise comparisons, identical across arms, plus two repeats
  (Section~\ref{sec:exclusions}) --- twelve screens in total.
\item Two seven-point items: perceived effort, and confidence that the system
  now understands their taste.
\item Debrief: the fitted taste profile and five recommendations.
\end{steps}

Total is about twelve minutes.

\subsection{Recruitment and sample size}

Recruitment via Prolific, restricted to fluent English speakers aged 18+, no
prior participation in a pilot of this study. Prolific rather than a university
pool: students are unrepresentative in exactly the dimension that matters here
(age strongly predicts film familiarity), and the pool would also make the
department's own students the subjects of their own coursework.

Power analysis for the primary test. Treating H1 as a two-sample comparison of
mean held-out log loss with a medium standardised effect $d = 0.5$, one-sided
$\alpha = .05$ and power $0.80$ requires
\[
  n \;=\; \frac{(z_{1-\alpha} + z_{1-\beta})^2 \cdot 2}{d^2}
    \;=\; \frac{(1.645 + 0.842)^2 \cdot 2}{0.25} \;\approx\; 49.5
\]
per arm, so \textbf{50 per arm}. Inflating by 20\,\% for the exclusions in
Section~\ref{sec:exclusions} and for abandonment --- higher in Design 2 under
H3 --- gives \textbf{150 recruited, 75 per arm}. $d = 0.5$ is chosen as the
smallest difference that would change what anyone builds; a smaller one would be
real but not actionable.

Payment is \pounds 2.00 for a twelve-minute task, above the Prolific minimum,
paid on completion and also on withdrawal.

\subsection{Exclusions and attention checks}
\label{sec:exclusions}

All criteria are fixed before collection begins.

\begin{steps}
\item \textbf{Test--retest.} Two of the ten held-out pairs are shown a second
  time with the films swapped left to right. Disagreeing on \emph{both} excludes
  the participant. The swap matters: a participant clicking the same side out of
  habit is caught, whereas an identical re-presentation would not catch them.
\item \textbf{Floor on effort.} Median task time below 1.5\,s in Design 1, or any
  ranking submitted unmodified from its presented order in Design 2.
\item \textbf{Non-completion.} The held-out block not finished. Counted and
  reported by arm, since H3 predicts the rate differs.
\end{steps}

Excluded participants are reported by arm and by reason. The rules are written
down in advance precisely so that they cannot be applied to whoever spoiled the
result.

\subsection{Analysis plan}

Pre-registered on OSF before the first participant.

\emph{Primary.} Welch's $t$-test on mean held-out log loss, Design 2 versus
Design 1, one-sided. Reported with a bootstrap 95\,\% CI on the difference; the
CI, not the $p$-value, is what a reader should act on. Mann--Whitney as a
sensitivity check if a Shapiro--Wilk test rejects normality.

\emph{H2.} A mixed-effects logistic model over the implied stage-wise choices in
Design 2, regressing correctness of the model's prediction on stage index $k$
with a random intercept per participant. A negative coefficient on $k$ is
position-dependent degradation. Equivalently and more directly: fit a separate
temperature per stage and test whether it increases in $k$.

\emph{H3.} Two-proportion test on abandonment; Mann--Whitney on the effort item.

\emph{H4.} Refit each participant's $\hat{w}$ on truncated prefixes of their
responses (2, 4, 6, 8 minutes) and compare held-out log-loss curves by arm. This
is why per-response timestamps are recorded rather than a single total.

Holm--Bonferroni across the three secondaries. The primary is not corrected: it
is a single pre-specified confirmatory test.

\emph{Reporting.} All four hypotheses are reported whichever way they come out,
along with exclusion counts by arm, and the anonymised response data.

\subsection{Pilot}

Ten participants, not pooled into the main analysis, to fix the two numbers the
design depends on: the median duration of a pairwise comparison and of a
ranking of ten. Those set the expected task counts in Table~2 and feed the
budget. The pilot also tests whether the instructions are understood without a
question, and whether the drag interaction works on a phone --- if it does not,
the study is desktop-only and that becomes a stated population restriction.

\subsection{Ethics and data protection}

Approval from the TUHH ethics committee before recruitment. Informed consent on
the first screen, with the withdrawal right stated in the same words used on the
study screen. No special-category data.

Recorded: arm, film indices, response order, per-task duration, viewing
frequency, two Likert items. Not recorded: name, email, IP address, Prolific ID
beyond the payment reconciliation held separately and destroyed once payment
clears. The dataset is pseudonymous under GDPR Art.~4(5); the lawful basis is
consent under Art.~6(1)(a), withdrawable, and withdrawal deletes the row.
Retention five years, then deletion; anonymised data published with the
pre-registration.

\subsection{Threats to validity}
\label{sec:threats}

\begin{terms}
\term{Unfamiliar films.} A participant who recognises neither film is guessing
  from metadata, which enters the model as noise attributed to $w$. Mitigated by
  the 25{,}000-vote floor. Not eliminated --- an ``I don't know either of these''
  option would eliminate it, at the cost of making the two arms structurally
  different, since there is no natural equivalent when ranking ten. The floor is
  the lesser compromise.
\term{Violations of the choice axiom.} Near-duplicate films in a set of ten
  split the vote in a way Plackett--Luce cannot represent, and this hits Design 2
  specifically --- so it is confounded with the manipulation. A post-hoc check
  correlates each set's internal feature similarity with that ranking's residual;
  if the effect is large, the finding is about ranking-with-random-sets rather
  than about ranking.
\term{Position-dependent noise.} Assuming it away would let Design 2 borrow
  confidence it has not earned. H2 measures it directly.
\term{Model misspecification, shared.} Utility is linear and additive, so
  ``I like horror but not long horror'' is inexpressible. This applies equally to
  both arms, so it threatens the generalisation to better preference models, not
  the internal comparison.
\term{Population.} Prolific participants are not a random sample of humanity.
  The claim is about these designs on this population, and is stated that way.
\term{Simulated, not observed.} The information figures and the simulation in
  Section~\ref{sec:sim} assume people behave as the model says. They set the
  direction of H1; they are not evidence for it. That is the whole reason for
  running the study.
\term{Timing assumptions.} The 6\,s and 90\,s per-task figures do the heavy
  lifting in every time-matched comparison here, and they are assumptions until
  the pilot measures them. If a ranking of ten takes 45\,s rather than 90\,s,
  Design 2 gets twice the tasks and the conditioning gap narrows sharply. The
  pilot exists mainly to pin these two numbers down, and the analysis is
  re-run against the measured values before recruitment.
\end{terms}

\subsection{What the simulation predicts}
\label{sec:sim}

Before recruiting anyone it is worth knowing whether the design can detect
anything, and in which direction. Simulated participants with a known $w$ answer
according to Plackett--Luce with a stage-dependent temperature $1 + \beta k$,
under the eight-minute budget at the assumed 6\,s and 90\,s per task. Both arms
are fitted by the same MAP procedure and scored on ten held-out pairs.

\begin{table}[H]
\centering\small
\begin{tabular}{lcccccc}
\toprule
& \multicolumn{2}{c}{\textbf{Held-out log loss}} & &
  \multicolumn{2}{c}{\textbf{Correlation with true $w$}} & \\
\cmidrule(lr){2-3}\cmidrule(lr){5-6}
\textbf{Position noise} & D1 & D2 & \textbf{Paired $d$} & D1 & D2 &
\textbf{Favours D2} \\
\midrule
$\beta = 0$ (none)      & 0.488 & 0.488 & $-0.00$ & 0.753 & 0.766 & 54\,\% \\
$\beta = 0.25$ (mild)   & 0.488 & 0.557 & $-0.34$ & 0.753 & 0.626 & 41\,\% \\
$\beta = 0.5$ (marked)  & 0.488 & 0.599 & $-0.51$ & 0.753 & 0.547 & 33\,\% \\
\bottomrule
\end{tabular}
\caption{120 simulated participants per row; 80 comparisons against 5 rankings
in the same eight minutes. Negative $d$ favours Design 1. $\beta = 0$ is the
model's own optimistic assumption --- every stage of a ranking answered as
carefully as the first. $\beta = 0.5$ makes the last stage roughly five times
noisier than the first.}
\end{table}

\paragraph{Reading it.} Even when the participant is assumed to rank ten films as
carefully as they would choose between two --- the most generous assumption
available, and one nobody believes --- Design 2 does not win. It ties, with a
paired $d$ of $-0.00$ and a coin-flip split across simulated participants. Its
higher correlation with the true $w$ at $\beta = 0$ (0.766 against 0.753) does
not convert into better held-out prediction, which is what the conditioning
result predicts: the extra accuracy sits in directions the held-out pairs rarely
probe, while the poorly-constrained directions cost it elsewhere.

Add any position-dependent carelessness and Design 1 pulls clear. At
$\beta = 0.25$ the paired $d$ is $-0.34$; at $\beta = 0.5$ it is $-0.51$, around
the effect the study is powered to detect.

\paragraph{What this does and does not license.} It sets the direction of H1, and
it says the study is worth running --- if Design 2 had won under every assumption,
the interesting question would have been settled at a desk. It is not evidence
about people. Every row assumes participants are Plackett--Luce agents with a
particular noise schedule, and $\beta$ for real participants is exactly the
unknown. What the simulation establishes is that the answer turns on $\beta$, and
$\beta$ can only be measured by asking people --- which is H2, and which is why
the interface records per-stage responses rather than just the final order.

% ============================================================================
\section{Task 4 --- The interface}

Implemented in \file{project4/}, reachable from the landing page at
\file{/project4/}, which offers this PDF on one side and the study on the other
as the brief requires.

\paragraph{Flow.} Consent and the frequency question
(\file{templates/project4/consent.html}) $\rightarrow$ blocked random assignment
$\rightarrow$ elicitation in the assigned design, \file{pair.html} or
\file{rank.html} $\rightarrow$ the twelve-screen held-out block, always
\file{pair.html} $\rightarrow$ debrief. The state machine is \file{study.py};
\file{next\_task} is the single place that decides what the participant sees
next, so the protocol can be read in one function.

\paragraph{Storage.} A \file{Participant} row and one \file{Response} row per
task, in the database rather than the session, because session-only storage
loses everything the moment a tab closes. A pairwise choice is stored as a
ranking of length two, so the schema has one representation and the fitting code
never branches on arm --- the practical payoff of Section~3's first
justification.

\paragraph{Reproducibility.} Every draw derives from one stored integer seed per
participant. The held-out block is therefore fixed for that participant:
reloading cannot reroll it into an easier one, and a whole session can be
reconstructed from the seed.

\paragraph{Accessibility.} Ranking supports drag-and-drop \emph{and} per-row
up/down buttons. This is not politeness. Drag-and-drop is unusable with a
keyboard or a screen reader, so a drag-only interface would exclude a subset of
participants --- and would do so invisibly, since the excluded participants
appear in the data as abandonments and inflate exactly the H3 statistic being
measured. The ranking form also serves its hidden inputs in the presented order,
so a participant with scripting disabled submits something coherent rather than
nothing.

\paragraph{Timing.} Per-task durations are measured server-side from when the
page was issued, capped at 600\,s so a participant who walks away does not
consume the whole budget in one task. H4 needs per-response timestamps, so they
are stored individually rather than accumulated.

\paragraph{Preview mode.} The consent screen carries a clearly separated block,
marked as not part of the participant's view, that forces an arm and shortens
the budget to 75 seconds so the interface can be walked end to end. Preview runs
are flagged on the participant row and can be filtered out of any analysis.

\paragraph{Explanations.} Twenty-four plain-language pages
(\file{project4/explain.py}), linked inline at the point each term first appears
rather than collected in a glossary nobody opens. The participant-facing pages
are covered too, deliberately: consent that requires a statistics degree to
understand is not informed consent.

% ============================================================================
\section{Limitations}

The linear utility model cannot express interactions between features, and the
representation cannot express anything about who made the film. Both are chosen
constraints, and both would be worth relaxing given a larger elicitation budget
than a new user will tolerate.

Uniform random sampling of films is what the brief permits and is not what a
deployed system should do; an adaptive scheme choosing maximally informative sets
would improve both arms. It is excluded here because it would interact with the
manipulation.

Most importantly, the study has not been run. Every quantitative claim about the
two designs in this report is a simulation under the model's own assumptions,
and the point of the design in Section~4 is that those assumptions are the thing
in doubt.

\end{document}
